% Generated by tools/build-kiwisdr-test-plan.py; edit the Markdown source.
% Structure adapted from acq-latex-templates engineering/verification/test-plan.
% Upstream revision fef160e37542fd7bc8e62de81404a83b8634a396.
\documentclass[11pt,letterpaper]{article}
\usepackage[T1]{fontenc}
\usepackage{lmodern,microtype}
\usepackage[margin=0.85in,includeheadfoot,headheight=20pt]{geometry}
\usepackage{booktabs,longtable,array,needspace}
\usepackage{enumitem,fancyhdr,lastpage,parskip,xcolor}
\usepackage{listings,tikz,graphicx,adjustbox,pdflscape,caption}
\usetikzlibrary{arrows.meta}
\PassOptionsToPackage{obeyspaces}{url}
\usepackage[hidelinks,pdftitle={Linkwise Modem KiwiSDR Test Plan},pdfauthor={Linkwise Project}]{hyperref}
\usepackage{xurl}
\definecolor{accent}{HTML}{216875}
\definecolor{pale}{HTML}{EFF5F6}
\setlength{\emergencystretch}{3em}
\setlength{\parindent}{0pt}
\setlength{\parskip}{5pt}
\setlist{leftmargin=*,itemsep=3pt,topsep=4pt}
\newlength{\tableavail}
\renewcommand{\arraystretch}{1.13}
\lstset{basicstyle=\ttfamily\small,breaklines=true,columns=fullflexible,
  backgroundcolor=\color{pale},frame=single,framerule=0pt,xleftmargin=4pt,xrightmargin=4pt}
\pagestyle{fancy}\fancyhf{}
\lhead{\small LINKWISE / MODEM TESTING}
\rhead{\small LW-MODEM-KIWI-001 / Draft 0.9}
\cfoot{\small\thepage\ of \pageref{LastPage}}
\setcounter{tocdepth}{2}

% All figures, diagrams and tables are landscape and bounded in both dimensions.
\captionsetup{font=small,skip=12pt}
\newcommand{\landscapegraphic}[3]{%
 \begin{landscape}\thispagestyle{plain}\centering
 \includegraphics[width=0.96\linewidth,height=0.76\textheight,keepaspectratio]{#1}\par
 \captionof{figure}{#2}\label{#3}
 \end{landscape}
}
\newcommand{\landscapetikz}[3]{%
 \begin{landscape}\thispagestyle{plain}\centering
 \begin{adjustbox}{max totalsize={0.96\linewidth}{0.76\textheight},center}
 \input{#1}
 \end{adjustbox}\par
 \captionof{figure}{#2}\label{#3}
 \end{landscape}
}
\begin{document}
\hypersetup{pageanchor=false}
\begin{titlepage}
\centering
\vspace*{1in}
{\Huge\bfseries Linkwise Modem\par}
\vspace{0.2in}
{\LARGE KiwiSDR Network Test Plan\par}
\vspace{0.5in}
{\large One-way RF observation and decoder regression\par}
\vfill
\begin{minipage}{0.7\textwidth}
\textbf{Document:} LW-MODEM-KIWI-001\par
\textbf{Revision:} 0.9 (Draft)\par
\textbf{Date:} 12 November 2026\par
\textbf{Prepared for:} Linkwise project\par
\textbf{Test lead / operator:} Not assigned\par
\textbf{Execution status:} Live receive evidence recorded. Controlled trial unscheduled.
\end{minipage}
\vspace{0.4in}
\begin{minipage}{0.88\textwidth}
Exploratory monitoring recovered one live BW500 frame from a KiwiSDR
recording. The controlled trial will measure exact recovery of known
transmissions. Bidirectional modem and Winlink qualification remain open.
\end{minipage}
\vspace{0.5in}
\end{titlepage}
\hypersetup{pageanchor=true}
\begin{landscape}\thispagestyle{plain}\centering
\begin{adjustbox}{max totalsize={0.96\linewidth}{0.76\textheight},center}
\begin{tabular}{@{}lll@{}}
\toprule Revision & Date & Change\\\midrule
0.1 & 9 November 2026 & Initial KiwiSDR experiment draft\\
0.2 & 9 November 2026 & Template selection and objective/resource tailoring\\
0.3 & 9 November 2026 & Template-structured LaTeX edition and rendered diagram\\
0.4 & 9 November 2026 & Landscape figures/tables; PlantUML and TikZ views\\
0.5 & 9 November 2026 & Julian-centered GIS map and candidate register\\
0.6 & 9 November 2026 & Exact command sequence, expected results and bench helper\\
0.7 & 9 November 2026 & Capture procedures and access checks\\
0.8 & 10 November 2026 & Live receive evidence and pilot acquisition regression\\
0.9 & 12 November 2026 & Bench evidence, PRBS9 pilot, and station readiness\\\bottomrule
\end{tabular}
\end{adjustbox}
\captionof{table}{Document revision history}
\end{landscape}

\begingroup\small
\tableofcontents
\endgroup
\clearpage

\section{Scope}
\label{tp:scope}

Use geographically separated KiwiSDR receivers to capture known Linkwise transmissions, measure one-way frame recovery under actual HF conditions, and build a reproducible decoder regression corpus. Test the native modem first; agent chat, global routing, and message-service behavior are outside this plan.

The controlled campaign scores a receiver observation when a known RF transmission overlaps a valid capture and Linkwise recovers the expected frame bytes. Exploratory monitoring verified one CRC-valid BW500 level-2 frame from the full Lake Merwin KiwiSDR recording on 7.105 MHz USB at about 5:25 p.m. PST on 10 November 2026. Six filter settings reproduced that same frame. The transmitter-side frame count and payload were unavailable, so this result establishes receive capability for one observed frame. The controlled trial will measure recovery rates with a known denominator. Two-way VARA and Winlink qualification require separate station tests. \href{../../../rust/docs/../../lab-notes/lab-notes.pdf}{Lab notes}

Use three distinct evidence labels:

\begin{itemize}
\item \textbf{Synthetic:} local waveform generation, impairment injection, and replay.
\item \textbf{RF observation:} remote recording of a physically transmitted known waveform.
\item \textbf{Operational link:} a separate, qualified two-station exchange, not established by this plan.
\end{itemize}
Record the radio-free mixed-peer bench separately from those campaign results. On merged main commit \path{70cf5ed}, BW2300 and BW2750 each passed three consecutive stock VARA HF 4.9.1 DATA recovery trials in both tested Linkwise roles, with exact bytes, two-ended disconnect, and a same-socket reconnect. Two BW500 Winlink Express sessions exchanged byte-exact 64-byte PRBS9 attachments, one session initiated by each modem. BW500 ACK-loss and mixed Pat repeat qualification on the merged build remain open. None of these bench rows enters the Kiwi RF denominator. \href{../../../rust/docs/../../docs/testing/vara-interoperability.md}{Wide recovery sign-off} and \href{../../../rust/docs/../../docs/media/modem/winlink-prbs9/verified-attachments.json}{Winlink PRBS9 evidence}.


\subsection{Implementation baseline and prerequisites}
\label{md:2}

The \href{../../../rust/docs/modem.md}{modem guide} is the implementation baseline. The current decoder accepts 48 kHz mono PCM16 or float32 WAVs. Levels 1--3 use a rolling pilot search with 0 and ±12 Hz frequency hypotheses. Level 4 uses the fixed 1\% RMS energy gate and requires an external measurement profile. It handles one frame per burst. The audio history is capped at 30 seconds, and the receiver does not track sample-clock drift. Preserve these limits in every report.

Available now:

\begin{itemize}
\item Native level 1--3 synthesis and streaming decode; optional level-4 components.
\item \path{linkwise decode-wav} and JSON frame reports containing level, control, payload hex, and sample ranges.
\item Local \href{../../../rust/docs/modem-demo.md}{SigPlot demo}, known waveform examples, and fixture tests.
\end{itemize}
Needed before an RF campaign:

\begin{itemize}
\item A qualified station playback/radio/PTT path, operator, and permitted test configuration. The web demo does not key a transmitter, and optional CPAL capture is not an integrated operational transmit service.
\item A cooperating receiver registry with owner-agreed automated capture windows, slot/time limits, settings, and retention/publication terms. Public listing alone is not permission for an unattended campaign.
\item A pinned KiwiClient revision, capture continuity instrumentation, deterministic conversion, manifest generation, and expected-frame comparison tooling.
\item A verified physical local transmit monitor recording, plus transmitter start, stop, and fault evidence. Generating a WAV alone does not confirm RF emission.
\end{itemize}
Do not automatically key a radio or occupy public receivers to execute this written plan. Campaign execution is a subsequent station operation.


\section{Referenced Documents}
\label{tp:references}

Primary technical sources reviewed on 9 November 2026:

\begin{itemize}
\item \href{https://github.com/jks-prv/kiwiclient/blob/master/README.md}{KiwiClient README}: stream/recording capabilities, resampling, and GNSS-tagged IQ format.
\item \href{https://github.com/jks-prv/kiwiclient/blob/master/kiwirecorder.py}{Kiwi recorder implementation}: effective gain/AGC, compression, resampling, and sample-writing behavior.
\item \href{../../../rust/docs/modem.md}{Linkwise modem implementation status} and the current Rust WAV CLI.
\end{itemize}
The source links follow upstream \path{master}; resolve and record an immutable revision before implementing the capture adapter. The supplied \href{http://kiwisdr.com/public/}{public directory} failed to load directly. Section 4.1 now records a dated mirror-derived shortlist cross-checked against receiver metadata; recording access and campaign readiness remain unconfirmed.


\section{Test Items}
\label{tp:items}

Freeze these items in each as-tested manifest. Experimental builds have no qualified production baseline; deviations from a future production system are not yet applicable.

\begin{landscape}\thispagestyle{plain}\centering
\setlength{\tabcolsep}{5pt}\renewcommand{\arraystretch}{1.25}
\setlength{\tableavail}{\dimexpr\linewidth-4\tabcolsep\relax}
\begin{adjustbox}{max totalsize={0.96\linewidth}{0.76\textheight},center}
\begin{tabular}{@{}>{\raggedright\arraybackslash}p{0.23\tableavail}>{\raggedright\arraybackslash}p{0.54\tableavail}>{\raggedright\arraybackslash}p{0.23\tableavail}@{}}
\toprule
\textbf{Item} & \textbf{Configuration identification} & \textbf{Status}\\
\midrule
Native modem & Git commit, build flags, waveform profile hash, frame format and expected bytes & Prototype under test\\[5pt]
Transmit chain & Source WAV, playback/radio settings, monitor recording and station logs & Physical execution path requires qualification\\[5pt]
Kiwi capture chain & Receiver endpoint/settings, client revision, sample continuity and timing evidence & Remote test instrumentation; not an ARQ peer\\[5pt]
Import/comparison harness & Demodulator/resampler versions, transforms, fixed decode profiles and oracle & New tooling pending K0\\[5pt]
\bottomrule
\end{tabular}
\end{adjustbox}
\captionof{table}{Test items and configuration evidence}
\end{landscape}

\landscapegraphic{../build/kiwisdr-test-plan/diagrams/Kiwi_Physical_Setup}{Proposed station setup: separate audio, RF, local monitoring and exclusive PTT paths. The monitor path remains to be qualified.}{fig:physical-setup}

\section{Test Approach}
\label{tp:approach}

\subsection{Measurement architecture}
\label{md:3}

Use KiwiClient/\path{kiwirecorder.py} as the initial capture adapter rather than browser audio playback or desktop sound recording. Upstream supports audio, IQ, and waterfall capture, including multiple receivers. IQ Kiwi-WAV recording can retain GNSS timestamp chunks; validity must be checked before timing use. The recorder exposes gain/AGC and compression settings and has multiple resampling paths. Pin the client and record which path actually executes. Sources: \href{https://github.com/jks-prv/kiwiclient/blob/master/README.md}{KiwiClient README}, \href{https://github.com/jks-prv/kiwiclient/blob/master/kiwirecorder.py}{recorder source}.

The publication edition includes five proposed PlantUML views (measurement chain, physical setup, capture-processing branches, observation scoring, and K0--K5 gates) and two TikZ schematics (burst timing and signal annotation). Sources are in \path{docs/testing/diagrams/} and \path{docs/testing/figures/}. The signal annotation is a conceptual illustration, not measured RF data; no receiver locations or global coverage map are invented before site selection. All views use landscape pages.

These are source capabilities, not evidence that any particular public receiver currently supports, permits, or has capacity for the selected capture profile.


\landscapegraphic{../build/kiwisdr-test-plan/diagrams/Kiwi_Measurement_Chain}{Proposed measurement chain. Kiwi receivers observe one-way RF; the expected-frame manifest is the independent oracle.}{fig:measurement-chain}

\subsection{Stimulus and campaign matrix}
\label{md:5}

Freeze payload bytes, control fields, waveform hashes, and schedule before RF. Each burst must have a unique expected frame identity. Encode a compact test sequence in the payload; the manifest holds full experiment metadata so it need not fit the smallest frame. Include the exact padding/control bytes in the oracle.

The pilot uses nine ASCII payload bytes before level-dependent padding. \path{K} and a two-digit index identify the burst. The remaining six bytes encode 30 successive PRBS9 bits as six five-bit Base32 symbols, most-significant bit first within each symbol. The right-shifting register starts at \path{0x1ff}, emits bit zero, and inserts the XOR of bits 7, 6, 4, and 0 at bit 8. The helper records each burst's starting register state and the waveform seed. The block follows a fixed balanced level order, so the manifest can reproduce the sequence without a randomized ordering seed. Retain the source hashes and calculated schedule with the station playback record.

Begin with one full DATA frame per burst at levels 1--3. Confirm every stimulus decodes locally before transmission. Maintain at least two seconds of quiet guard between bursts and record background before and after the block. Bound individual test bursts below 20 seconds and preserve enough guard audio to close the detector's candidate; do not decode an arbitrarily truncated burst.

\begin{landscape}\thispagestyle{plain}\centering
\setlength{\tabcolsep}{5pt}\renewcommand{\arraystretch}{1.25}
\setlength{\tableavail}{\dimexpr\linewidth-6\tabcolsep\relax}
\begin{adjustbox}{max totalsize={0.96\linewidth}{0.76\textheight},center}
\begin{tabular}{@{}>{\raggedright\arraybackslash}p{0.18\tableavail}>{\raggedright\arraybackslash}p{0.25\tableavail}>{\raggedright\arraybackslash}p{0.32\tableavail}>{\raggedright\arraybackslash}p{0.25\tableavail}@{}}
\toprule
\textbf{Phase} & \textbf{Stimulus / comparison} & \textbf{Repetitions and scope} & \textbf{Purpose}\\
\midrule
K0: offline qualification & Known clean waveforms, round-trip conversion, sample gap controls, off-target/background controls & Levels 1--3; no RF or public slots required & Qualify tooling and the exact-byte oracle\\[5pt]
K1: receiver preflight & Background plus one scheduled known level-1 burst & One cooperating receiver and one valid observation window & Verify tuning, capture continuity, and acquisition scale\\[5pt]
K2: pilot & Interleaved levels 1--3, fixed station power/configuration & 10 bursts per level, up to three simultaneous receivers & Measure one-way recovery against known transmissions\\[5pt]
K3: temporal replication & Repeat K2 in distinct contact windows & At least three sessions; record dates and UTC separately & Expose time/path variation without pooling it away\\[5pt]
K4: decoder comparison & Identical retained K2/K3 captures replayed through baseline and candidate commits & All valid captures; paired comparison & Isolate software changes from HF variation\\[5pt]
K5: diagnostic extension & One variable at a time: bounded offset, power setting, AGC profile, BW500 level 4, or wide bandwidth & Separate blocks after K0--K2, with a mode-specific local oracle and verified receiver passband & Characterize specific failure modes\\[5pt]
\bottomrule
\end{tabular}
\end{adjustbox}
\captionof{table}{Campaign matrix and experimental purpose}
\end{landscape}
\landscapetikz{figures/burst-timeline}{Illustrative timing: continuous pre-roll and post-roll, uniquely identified bursts and quiet guards. Spacing is schematic, not proportional to seconds.}{fig:burst-timeline}

A complete K2 block contains 30 transmitted bursts and up to 90 receiver/burst observations, not 90 independent transmissions. K3 repeats that block across sessions. Randomize or balance level order within each block and preserve the ordering seed. Final block length and airtime are computed from actual waveform durations before scheduling; abort if they exceed station or receiver limits.

Do not vary power, passband, gain, level, and software simultaneously and then attribute the result to one variable. Native transmit-version A/B comparisons need interleaved, balanced RF trials; unlike decoder replay, different RF transmissions do not experience identical channels.


\subsection{Per-session execution procedure}
\label{md:6}

\begin{enumerate}
\item Assign experiment, session, burst, and receiver IDs; freeze source commits, configuration, expected bytes, hashes, and the scheduled listening window.
\item Run existing native tests and decode every source WAV. Pin the external profile if level 4 is included. Preserve machine-readable outputs.
\item Preflight receiver access and effective settings. Start capture at least 30 seconds before the first scheduled transmission, without squelch. Confirm samples arrive and establish a noise/background interval.
\item Confirm the qualified station configuration and local monitor. Transmit the bounded schedule through the station's approved execution path, retaining actual playback/RF evidence and recording aborted or partial transmissions.
\item Continue capture at least 30 seconds after the last burst; terminate and release receiver slots at the registered deadline.
\item Hash originals, validate continuity/settings, and classify each observation before looking at decoder success. Preserve all exclusions with reasons.
\item Derive 48 kHz mono analysis WAVs with a pinned high-quality resampler and, for IQ, a pinned demodulation pipeline. Record delay and original sample/time mapping. Keep each eligible continuous interval intact, including background and guards. Split at actual discontinuities and retain their source ranges.
\item Decode both baseline and candidate versions and compare exact frame bytes. Generate the report, archive required artifacts, and retain failures as well as successful examples. Publish only material covered by the capture terms.
\end{enumerate}
\subsubsection{Station readiness record before K1}
The test lead records the disposition, station operator, capture operator, procedure revision, date and UTC window. The station operator confirms the as-tested radio, firmware, audio devices and gain, route, antenna or load, mode, frequency, and power. Use the station's permitted operating envelope.

Before the first radiated block, rehearse the playback and abort sequence with transmit inhibited or into a suitable conducted load. Record the exact source hash, local monitor trace, PTT assertion and release, measured turnaround, and the response to an operator stop. Record the limit and observed value for each timing check. A failed release, unexpected transmission, monitor loss, wrong route or frequency, or receiver-window loss stops the block. The operator confirms PTT release and restores the station before dispositioning the attempt.

Record each cooperating receiver's agreed slot, time limit, passband, gain profile, and retention or publication terms. Recheck access and effective settings near the scheduled window. The station operator authorizes the bounded K1 burst only after the test lead has recorded K0, station readiness, and receiver readiness. Authorize K2 after reviewing K1 capture validity and the disposition of any decode miss.

No experiment starts solely because a model predicts an opening. The station execution path retains local authority, timing, and stop conditions.


\subsection{Capture validity and normalization}
\label{md:7}

Assign validity independently for decode, amplitude/SNR, and precise timing. For example, a contiguous audio capture with uncertain clock can support frame recovery while remaining ineligible for propagation timing measurements.

\begin{itemize}
\item Sequence discontinuities, reconnects, missing samples, unexplained rate changes, and truncated windows make affected intervals decode-ineligible unless complete unaffected bursts can be isolated with explicit provenance. Do not interpolate missing signal samples and count the result as a valid RF observation.
\item For the offline continuity check, pair a clean local fixture with a copy missing a 50 ms interior block. At 12 kHz this is 600 samples. Preserve original sample indices or packet sequence metadata and the deletion range. Require the continuity check to flag the affected burst before decoding. A wide-mode stock turn-grant capture lost 2,400 samples at 48 kHz and timed out, motivating this control. WAV length and decoder success alone cannot locate an internal gap in a live recording.
\item Receiver unavailable/busy, permission denial, client faults, and inability to verify required settings are infrastructure outcomes. A valid continuous recording with no recovered signal remains an RF/decode outcome.
\item Preserve overloaded/clipped captures and their flags. Exclude them from clean sensitivity claims under predeclared rules; report their prevalence separately.
\item Never use internet packet arrival time as RF arrival time. GNSS validity, sample continuity, and calibration determine timing eligibility. This pilot does not attempt TDoA or absolute propagation-delay measurement.
\end{itemize}
Run two explicitly different decoder profiles on each eligible capture:

\begin{enumerate}
\item \textbf{As received:} fixed conversion scale; no per-burst amplitude normalization, denoising, carrier correction, or timing correction beyond the declared import. This exposes acquisition behavior on the original capture scale.
\item \textbf{Diagnostic:} optional documented gain, frequency, or sample-rate correction. Save the transform and never mix these results into the as-received score.
\end{enumerate}
Converting an actual lower-rate recording to 48 kHz changes representation; it does not restore missing bandwidth or eliminate oscillator/sample-clock errors.

\subsubsection{Regression replays before the next RF session}
Build Linkwise once and replay each full 48 kHz mono input using the existing \path{linkwise decode-wav} command. Preserve the IQ demodulation settings and resample at the recorded input rate before importing IQ. The decoder finds level-1--3 boundaries and its own frequency hypotheses. The regression supplies no frame windows or manual frequency correction.

\begin{itemize}
\item \textbf{Lake Merwin live capture.} Expect exactly one level-2 frame from the full 7.105 MHz USB derivative. Check the 23-byte body hash in the evidence record. The IQ center is 7.1065 MHz. Both 330 Hz and 1200 Hz baseband low-pass derivatives are verified profiles.
\item \textbf{Published off-air exchange.} Replay the full slot 13 recording. Require three frames in order, levels 2, 3, 2, with the recorded body hashes.
\item \textbf{Kiwi noise controls.} Require byte-exact recovery of each +6 dB stock level-1 and level-2 injection. Require zero frames from the unmodified Maui noise capture.
\end{itemize}
Compare SHA-256 over decoded payload-field bytes followed by the one control byte. Check count and order as well as hashes. Report lost or added frames, duplicated frames, and truncated energy candidates. These are regression cases, separate from the frozen holdout and the live recovery denominator. The archived automatic acquisition JSON and preprocessing evidence are in the local follow-up directory listed in the lab notes.


\landscapegraphic{../build/kiwisdr-test-plan/diagrams/Kiwi_Capture_Processing}{Proposed processing branches. Preserve originals and continuity evidence; score as-received and diagnostic outputs separately.}{fig:capture-processing}

\begin{landscape}\thispagestyle{plain}\centering
\subsection{Metrics and scoring}
\label{md:8}

\setlength{\tabcolsep}{5pt}\renewcommand{\arraystretch}{1.25}
\setlength{\tableavail}{\dimexpr\linewidth-4\tabcolsep\relax}
\begin{adjustbox}{max totalsize={0.96\linewidth}{0.76\textheight},center}
\begin{tabular}{@{}>{\raggedright\arraybackslash}p{0.23\tableavail}>{\raggedright\arraybackslash}p{0.54\tableavail}>{\raggedright\arraybackslash}p{0.23\tableavail}@{}}
\toprule
\textbf{Metric} & \textbf{Definition and denominator} & \textbf{Availability}\\
\midrule
Capture availability & Complete eligible receiver/burst windows / planned windows & New capture manifest and validity checker\\[5pt]
Exact frame recovery & Unique exact expected frames / expected frames confirmed transmitted within decode-eligible windows & Current decoder output plus new comparator\\[5pt]
Acquisition yield & Expected bursts associated with receiver candidates / eligible confirmed transmitted bursts & Requires additional candidate/frame association reporting\\[5pt]
Nonmatching CRC-valid output & Valid decoded frames conflicting with the expected identity/bytes & Comparator; retain sample positions and investigate every case\\[5pt]
Background false detections & CRC-valid frames during confirmed no-test-transmission intervals, per recorded hour & Current decoder plus background labels; unrelated traffic is classified separately\\[5pt]
Offline decode cost & Wall/CPU time and peak memory per audio duration, fixed machine/configuration & New harness; exclude RF pacing/network wait\\[5pt]
Frequency offset / fading & Estimated against the known signal, with method and uncertainty & Later instrumentation; not currently exported as qualified metrics\\[5pt]
SNR & Declared in-band signal/noise estimate with bandwidth, method, and uncertainty & Separate estimator; not a raw S-meter comparison\\[5pt]
\bottomrule
\end{tabular}
\end{adjustbox}
\captionof{table}{Metrics, denominators and implementation availability}
\end{landscape}
Report the number of confirmed complete transmissions, eligible captured observations, and unique exact recoveries separately.

Missing expected frames count as failures when capture and transmission evidence are valid, even if acquisition finds nothing. Count a repeated decoded frame only once for recovery; separately count duplicates. A partial transmission is reported as such, not included as an expected complete frame. Retain orphan frames and unknown associations rather than silently forcing matches.

Do not claim BER from CRC-valid frames alone. Failed frames do not expose all bit errors through the existing API. Do not label unique recovered bytes per transmitted airtime as ARQ goodput or end-to-end message throughput.

For SNR, use a declared signal band and comparable noise intervals with the same receiver settings and bandwidth. Where justified, estimate signal power as on-burst power minus estimated noise power in linear units. If interference, AGC changes, nonstationary noise, or nonpositive subtraction invalidate the estimate, report SNR unknown. Display SNR bandwidth in every comparison. Do not infer calibrated path loss from different receivers' S-meter readings.

Report per receiver, level, band, UTC session, capture profile, and software revision. Include raw counts and paired baseline/candidate changes. Repeated frames on one fading path and simultaneous observations of one transmission are correlated; use session/contact-block summaries or a justified clustered interval method. A pilot's high score is not a global reliability estimate.


\landscapegraphic{../build/kiwisdr-test-plan/diagrams/Kiwi_Observation_Scoring}{Observation classification before scoring: confirmed transmission and eligible capture determine the denominator; a valid no-decode window remains a miss.}{fig:observation-scoring}

\landscapetikz{figures/signal-annotation}{Conceptual annotation example, not measured RF data. Burst windows and noise references must be independently defined; visual energy alone is not exact recovery.}{fig:signal-annotation}

\subsection{Acceptance and progression gates}
\label{md:9}

These are proposed experimental gates, not achieved modem service levels. Pre-register any change before collecting the evaluation dataset.

\begin{itemize}
\item \textbf{K0 tooling gate:} all clean level 1--3 fixtures survive supported import and conversion with identical expected bytes; repeated replay produces identical frame results; seeded capture gaps are identified and not silently repaired; metadata/hashes and original-to-derived sample mapping are complete. The full-recording regression cases in section 7 must also meet their exact frame-count, order, and hash expectations.
\item \textbf{K1 capture gate:} one correctly tuned, contiguous observation includes its full known RF burst and guards, with local transmission confirmation and usable settings evidence. A failed decode can still be a valid diagnostic observation; it is not a receiver preflight failure solely for decoding poorly.
\item \textbf{K2 pilot completion:} all planned observations have an outcome, exclusion reason, or explicit missing status; at least one exact RF-observed frame is needed before claiming remote RF decoding demonstrated. Missing that result triggers diagnosis, not removal of failed windows from the denominator.
\item \textbf{K3 replication:} collect three distinct session blocks with complete provenance before describing repeatable behavior for the sampled paths. No universal minimum frame-recovery percentage is asserted without a defined channel and station operating envelope.
\item \textbf{K4 release comparison:} on a frozen holdout corpus, no new exact-byte mismatch, unreported discontinuity, crash, or out-of-budget resource use. Report all gained and lost frames. A recovery-loss tradeoff needs an explicit profile-specific decision; aggregate improvement cannot hide a lost case.
\item \textbf{K5 extension:} level 4, DSP corrections, and transmit comparisons remain separately labeled. Never promote synthetic-only evidence to RF or stock-peer qualification.
\end{itemize}
If sensitivity targets are desired, define them after the pilot using a separate training corpus, then freeze thresholds and test on held-out captures. Do not select the easiest receivers or adjust gain/thresholds on evaluation recordings and report the outcome as unbiased performance.


\landscapegraphic{../build/kiwisdr-test-plan/diagrams/Kiwi_Qualification_Gates}{Proposed progression from tooling through pilot, replication, replay and optional diagnostics. These gates do not qualify a bidirectional operational link.}{fig:qualification-gates}

\begin{landscape}\thispagestyle{plain}\centering
\subsection{Objective traceability}
\label{tp:objectives}
These IDs are local experimental objectives, not inherited software requirements or the agent-chat requirements in the whitepaper. Method codes: T = test, A = analysis, I = inspection. Gate IDs retain the definitions in Section~\ref{md:9}.


\setlength{\tabcolsep}{5pt}\renewcommand{\arraystretch}{1.25}
\setlength{\tableavail}{\dimexpr\linewidth-6\tabcolsep\relax}
\begin{adjustbox}{max totalsize={0.96\linewidth}{0.76\textheight},center}
\begin{tabular}{@{}>{\raggedright\arraybackslash}p{0.18\tableavail}>{\raggedright\arraybackslash}p{0.25\tableavail}>{\raggedright\arraybackslash}p{0.32\tableavail}>{\raggedright\arraybackslash}p{0.25\tableavail}@{}}
\toprule
\textbf{Objective} & \textbf{Method / phase} & \textbf{Entry and exit criterion} & \textbf{Required evidence}\\
\midrule
KW-OBJ-01: trustworthy acquisition and import & I/T, K0--K1 & Enter with pinned tooling and known fixtures; exit with K0 tooling criteria met and K1 capture validity established & Hashes, versions, effective settings, continuity results, conversion/sample mapping\\[5pt]
KW-OBJ-02: exact one-way RF frame recovery & T/A, K2 & Enter after tooling/capture readiness; complete the planned matrix and report every observation; claim RF recovery only after an exact observed frame & Transmit confirmation, expected bytes, capture eligibility, exact matches/misses and exclusions\\[5pt]
KW-OBJ-03: sampled path/time repeatability & T/A, K3 & Enter with K2 evidence; exit with three distinct complete session blocks & Receiver/path/session-stratified counts and uncertainty; no global reliability inference\\[5pt]
KW-OBJ-04: decoder regression comparison & T/A, K4 & Enter with a frozen holdout corpus and fixed machine/profile; apply K4 release criteria & Paired gained/lost frames, mismatches, failures, resource costs, and disposition\\[5pt]
KW-OBJ-05: diagnostic characterization & T/A, K5 & Enter after K0--K2 with a predeclared single-variable matrix; complete it or record why it stopped & Raw/corrected separation, estimator validity, all variable settings and unresolved limits\\[5pt]
\bottomrule
\end{tabular}
\end{adjustbox}
\captionof{table}{Objective traceability and required evidence}
\end{landscape}
A test result references the objective, phase/gate, plan revision, procedure revision, configuration manifest, and raw/reduced evidence. Deviations receive unique IDs, rationale, decision owner, and impact on validity; never discard a failed observation merely because a later repeat succeeds.


\section{Resources, Facilities, and Equipment}
\label{tp:resources}

\subsection{Receiver selection and recording profiles}
\label{md:4}

For the first pilot, seek three cooperating receivers: one relatively short path and two geographically different longer paths where available. At least one intercontinental observation is a later coverage objective, not a prerequisite for validating the capture pipeline. Record distance, bearing, UTC window, antenna information, and path evidence; do not infer availability from distance.

Select receivers using a registry snapshot, then preflight each endpoint. Do not hard-code an assumed global receiver count. Re-check access rules, frequency coverage, time limits, slots, clock information, and effective settings at each session. A failed preflight is an infrastructure exclusion, not an RF non-detection. Cap the pilot at three simultaneous streams, one slot per receiver, with an explicit stop deadline; never evade passwords, kicks, or time limits.

\textbf{Profile A --- initial audio capture:} uncompressed USB audio; no squelch; noise blanking, noise reduction, de-emphasis and other optional processing off where supported. Use a fixed documented passband that contains the complete synthesized waveform, with margin. Fix gain where practical; otherwise record AGC configuration and label the run accordingly. Verify gain does not clip.

\textbf{Profile B --- later diagnostic IQ capture:} original complex samples with recorded center frequency, actual sample rate, IQ convention, and available GNSS metadata. Offline demodulation must document frequency translation, sideband, filters, scaling, and delay. Preserve custom timestamp chunks in the original; do not assume a stereo IQ WAV is directly decodable by Linkwise.

Record the carrier/dial frequency and the waveform's audio-frequency range separately. Validate the RF-to-audio mapping before judging acquisition; do not confuse a receiver's passband-center setting with its USB carrier setting.

\subsubsection{Julian transmitter and candidate receiver map}
The transmitting station will be in \textbf{Julian, California}. The map uses 33.0786552° N, 116.6019637° W as an approximate Julian locality reference, not the actual antenna location. Replace it with the station coordinates before final path analysis or campaign execution. The locality reference is \href{https://www.topozone.com/california/san-diego-ca/city/julian/}{reported from USGS GNIS}.

\landscapegraphic{figures/kiwisdr-candidate-map.pdf}{GIS candidate network centered on approximate Julian, California. Directory-derived sites were cross-checked through metadata only; geodesics are not propagation forecasts or confirmed links.}{fig:candidate-map}
In the \href{../../../rust/docs/../../docs/testing/figures/kiwisdr-candidate-map.pdf}{interactive PDF map}, click a receiver ID or its candidate-register name to open its web interface. The full LaTeX/PDF test plan preserves these links; the PNG preview is static.

Accessibility recheck, \textbf{9 November 2026, 23:54 UTC}: all \textbf{11/11} candidates returned HTTP 200 for both their web interface and \path{/status}; each web response contained the Kiwi interface and every status response reported \path{active} with \path{offline=no}. The \href{../../../rust/docs/../../docs/testing/data/kiwisdr-access-check.json}{timestamped per-receiver results} include effective URLs, response codes and returned metadata. Reproduce from the repository root with \path{python3 tools/check-kiwisdr-candidates.py}; it exits nonzero if any candidate fails either check. These bounded HTTP requests open no audio streams and establish reachability only at the recorded time. Repeat before a capture session; access rules and audio-slot readiness still apply.

The map is a GIS product generated with GeoPandas, Cartopy and PROJ, using Natural Earth geographic data. Its global view is WGS84 azimuthal equidistant, centered on Julian, with a western-US inset. Distances and initial true bearings are shortest WGS84 ellipsoidal geodesics. Lines show candidate geometry, not predicted propagation, reliable contact, RF hop counts, or measured coverage. South Africa is a stretch objective, outside the initial pilot denominator.

Cross-reference snapshot: \textbf{9 November 2026, 23:28:19 UTC}, from the public \href{http://rx.linkfanel.net/kiwisdr_com.js}{KiwiSDR directory mirror}. The mirror header identifies the \href{http://kiwisdr.com/public/}{official public directory} as its source. That official endpoint still did not respond directly during this selection. All eleven chosen records were listed active; bounded metadata-only \path{/status} requests also returned active status with matching names and coordinates. No audio connections or receiver reservations were opened. This does not establish owner permission, future slot availability, clock quality, antenna calibration, or decodability from Julian. Coordinates remain owner-reported.

The initial pilot pool is \textbf{R02 KFS Omni, R03 Northern Utah, and R05 Maui}, subject to readiness and access arrangements. R01 San Diego is a nearby diagnostic control; terrain and low-angle/skywave behavior are not assumed. R04 and R06--R10 extend the sampled geography later. Co-located receiver channels and the Utah IP-address alias are not counted as independent path sites.

\begin{landscape}\thispagestyle{plain}\centering
\setlength{\tabcolsep}{5pt}\renewcommand{\arraystretch}{1.25}
\setlength{\tableavail}{\dimexpr\linewidth-6\tabcolsep\relax}
\begin{adjustbox}{max totalsize={0.96\linewidth}{0.76\textheight},center}
\begin{tabular}{@{}>{\raggedright\arraybackslash}p{0.12\tableavail}>{\raggedright\arraybackslash}p{0.4\tableavail}>{\raggedright\arraybackslash}p{0.23\tableavail}>{\raggedright\arraybackslash}p{0.25\tableavail}@{}}
\toprule
\textbf{Candidate} & \textbf{Site / experiment role} & \textbf{Distance / initial bearing} & \textbf{Receiver endpoint}\\
\midrule
R01 & San Diego, CA / control & 60 km / 246° true & \href{http://kiwisdr.ka9q.net:8073}{Open receiver}\\[5pt]
R02 & Half Moon Bay, CA / KFS Omni / pilot & 713 km / 314° true & \href{http://kiwisdr.kfsdr.com:8073}{Open receiver}\\[5pt]
R03 & Northern Utah / Kiwi 1 / pilot & 1,019 km / 21° true & \href{http://kiwisdr1.sdrutah.org:8073}{Open receiver}\\[5pt]
R04 & Winthrop, WA / MVARC / expanded & 1,741 km / 351° true & \href{http://mvarc.proxy.kiwisdr.com}{Open receiver}\\[5pt]
R05 & Maui, Hawaii / pilot & 4,147 km / 261° true & \href{http://22215.proxy.kiwisdr.com}{Open receiver}\\[5pt]
R06 & Washington DC area / NA5B / expanded & 3,582 km / 68° true & \href{http://www.na5b.com:8073}{Open receiver}\\[5pt]
R07 & Tokyo, Japan / JK1LOT / expanded & 9,048 km / 307° true & \href{http://jk1lot.ham-radio-op.net:8073}{Open receiver}\\[5pt]
R08 & Marahau, New Zealand / Kiwi 1 / expanded & 10,921 km / 226° true & \href{http://kiwisdr1.owdjim.gen.nz:8073}{Open receiver}\\[5pt]
R09 & Santiago, Chile / expanded & 8,796 km / 142° true & \href{http://ce3pbr.proxy.kiwisdr.com}{Open receiver}\\[5pt]
R10 & Wessex, United Kingdom / G8JNJ / expanded & 8,698 km / 36° true & \href{http://wessex.zapto.org:8073}{Open receiver}\\[5pt]
R11 & Gqeberha, South Africa / ZS2JB / stretch & 16,538 km / 102° true & \href{http://alg.twrmon.net:8073}{Open receiver}\\[5pt]
\bottomrule
\end{tabular}
\end{adjustbox}
\captionof{table}{Julian-centered receiver shortlist; metadata checked, access unconfirmed}
\end{landscape}
The \href{../../../rust/docs/../../docs/testing/data/kiwisdr-candidates.json}{candidate registry} retains stable listing IDs, original names, endpoints, owner-reported coordinates, listing timestamps, metadata-check evidence, selection rationale, and the source snapshot SHA-256. The companion CSV, GeoJSON, and GeoPackage can be inspected in QGIS; see the \href{../../../rust/docs/../../docs/testing/data/README.md}{GIS provenance guide}. The full raw directory snapshot is local build evidence, not a newly asserted worldwide inventory count. Refresh registry and endpoint checks before execution.


\subsection{Resource and personnel status}

\begin{landscape}\thispagestyle{plain}\centering
\setlength{\tabcolsep}{5pt}\renewcommand{\arraystretch}{1.25}
\setlength{\tableavail}{\dimexpr\linewidth-6\tabcolsep\relax}
\begin{adjustbox}{max totalsize={0.96\linewidth}{0.76\textheight},center}
\begin{tabular}{@{}>{\raggedright\arraybackslash}p{0.24\tableavail}>{\raggedright\arraybackslash}p{0.12\tableavail}>{\raggedright\arraybackslash}p{0.36\tableavail}>{\raggedright\arraybackslash}p{0.28\tableavail}@{}}
\toprule
\textbf{Resource / role} & \textbf{Pilot quantity} & \textbf{Required responsibility / evidence} & \textbf{Current status}\\
\midrule
Qualified transmit station and operator & 1 & Approved configuration, playback/RF monitor, PTT release and stop procedure & Not assigned or qualified under this plan\\[5pt]
Cooperating Kiwi receivers / owners & Up to 3 & Agreed windows, access/automation limits, effective settings and capture terms & Shortlisted in section 4.1; windows/access unconfirmed\\[5pt]
Test lead & 1 role & Freeze matrix, review readiness, record deviations and terminal outcomes & Not assigned\\[5pt]
Capture/analysis operator & 1 role & Bounded recording, immutable artifacts, conversion and scoring & Not assigned; may share personnel with test lead\\[5pt]
Analysis host and storage & 1 environment & Versioned dependencies, adequate disk/network capacity, measured resource limits & Execution baseline not frozen\\[5pt]
Reference/calibration evidence & Per measurement & Clock/gain/frequency confidence appropriate to claimed metric & To be established; uncalibrated metrics stay excluded/unknown\\[5pt]
\bottomrule
\end{tabular}
\end{adjustbox}
\captionof{table}{Pilot resources, responsibilities and current readiness}
\end{landscape}

\section{Safety}
\label{tp:safety}

Abort the affected block on station safety/PTT faults, exceeded airtime or capture deadline, receiver-owner termination, or storage/continuity failures that prevent valid observation. The station's local stop/release procedure remains authoritative. Record partial transmissions, stop capture jobs, release receiver slots, preserve logs, and verify station restoration. Decide whether a fresh block is warranted only after documenting the cause and revised readiness; do not silently extend receiver sessions or repeat RF transmissions.


\section{Schedule}
\label{tp:schedule}

No calendar dates or receiver windows are committed. This schedule is dependency-based; each event starts only when its entry criteria are met. The pilot remains unscheduled.

\begin{landscape}\thispagestyle{plain}\centering
\setlength{\tabcolsep}{5pt}\renewcommand{\arraystretch}{1.25}
\setlength{\tableavail}{\dimexpr\linewidth-4\tabcolsep\relax}
\begin{adjustbox}{max totalsize={0.96\linewidth}{0.76\textheight},center}
\begin{tabular}{@{}>{\raggedright\arraybackslash}p{0.23\tableavail}>{\raggedright\arraybackslash}p{0.54\tableavail}>{\raggedright\arraybackslash}p{0.23\tableavail}@{}}
\toprule
\textbf{Event} & \textbf{Entry condition} & \textbf{Exit / next event}\\
\midrule
Tooling preparation and K0 & Frozen stimuli and pinned capture/import tools & K0 tooling evidence accepted\\[5pt]
Campaign readiness review & K0 complete; resources, windows, station operation and limits confirmed & Registered K1 window\\[5pt]
K1 and K2 capture & Valid preflight, qualified station chain and readiness review & Pilot outcomes and anomaly records\\[5pt]
K3 replication & K2 dataset and reviewed disposition & Three complete session blocks\\[5pt]
K4 replay and report & Frozen holdout corpus; baseline/candidate builds & Paired results, exclusions and disposition\\[5pt]
Optional K5 block & Separate variable matrix and readiness & Diagnostic report, no stronger qualification claim\\[5pt]
\bottomrule
\end{tabular}
\end{adjustbox}
\captionof{table}{Dependency-based schedule; no calendar windows committed}
\end{landscape}

\subsection{Readiness review}
\label{tp:readiness}

The readiness review must establish all of the following before a scheduled RF block: K0 tooling complete; test lead/operator assigned; station configuration and local monitor qualified; receiver windows confirmed; expected-frame matrix and airtime limits frozen; storage capacity checked; stop/restoration actions executable; and capture terms and measurement limitations recorded. The review is a campaign preparation step, not a new per-message approval mechanism.


\section{Data Management and Reporting}
\label{md:10}

Store large originals outside the source repository until storage and redistribution terms are settled. Suggested logical structure:

\begin{lstlisting}
experiment-id/
  plan.json                   # revision, matrix, limits, selection rules
  transmit/                   # expected frames, waveform, local monitor, logs
  receiver-id/session-id/
    capture.json              # effective settings, continuity, timing validity
    original/                 # immutable audio/IQ and sidecar logs
    derived/                  # 48 kHz WAVs and exact conversion recipe
    results/baseline.json
    results/candidate.json
  report.json                 # denominators, exclusions, comparisons, uncertainty
  report.md                   # interpretation and bounded claims
\end{lstlisting}
Each manifest needs commit IDs, dependency/profile versions, SHA-256 hashes, planned and actual UTC times with clock confidence, transmitter configuration, receiver endpoint/location provenance, mode/passband/gain settings, source sample rate, sequence-gap ranges, processing transforms, expected-frame IDs/bytes, transmission evidence, observation eligibility, and terminal reasons. Keep credentials outside exported manifests.

The report must show infrastructure availability beside RF recovery, list every invalid/missing observation, distinguish raw and corrected results, and retain receiver/path/session strata. Include waterfall/spectrum views as explanation, not as substitutes for exact-byte results.


For execution, develop a separately revisioned procedure using the repository's \href{https://github.com/8r4n/acq-latex-templates/blob/fef160e37542fd7bc8e62de81404a83b8634a396/templates/engineering/verification/test-procedure/main.tex}{Test Procedure template}, with step, action, expected result, actual result, and disposition fields. Report observations through the \href{https://github.com/8r4n/acq-latex-templates/blob/fef160e37542fd7bc8e62de81404a83b8634a396/templates/engineering/verification/test-report/main.tex}{Test Report template}, including as-tested configuration, deviations, anomalies, and evidence-backed conclusions. These companion deliverables are planned, not yet executable or completed records. A program-level TEMP is unnecessary for this bounded pilot; its broader strategy role remains with the whitepaper and implementation roadmap.


\section{Exact Test Sequence and Expected Results}
\label{tp:procedure}

These are the as-planned steps, not completed test records. Run commands from the repository root in a POSIX shell. Stop on any unexpected exit status; retain stdout, stderr, actual result, UTC observation time, operator, and disposition against the step ID. Successful process exit alone is not RF or delivery success. All command tables are landscape in the publication edition.

\textbf{Execution boundary:} A and B are local/file operations. C installs and checks a pinned recorder without opening a receiver. D/E require an agreed receiver window, qualified station procedure, and completed readiness review; they have not been executed. F requires valid acquired evidence. No command below keys a transmitter.

The test lead must provide actual antenna coordinates, permitted frequency, receiver passband/gain, access windows, and calibrated station procedure before RF. Commands with shell variables are exact invocations after those values are reviewed. Empty or unreviewed parameters are a stop condition, not permission to substitute a convenient frequency. Hardware-specific playback/PTT commands cannot be prescribed until that station interface is qualified.

\subsubsection{A --- Freeze a local baseline}
Use a new experiment directory for each attempt; do not reuse the example after it contains evidence. Terminal A is reserved for the local demo service. Terminal B runs the remaining shell commands and keeps the exported variables.

\begin{landscape}\thispagestyle{plain}\centering
\setlength{\tabcolsep}{5pt}\renewcommand{\arraystretch}{1.25}
\setlength{\tableavail}{\dimexpr\linewidth-6\tabcolsep\relax}
\begin{adjustbox}{max totalsize={0.96\linewidth}{0.76\textheight},center}
\begin{tabular}{@{}>{\raggedright\arraybackslash}p{0.1\tableavail}>{\raggedright\arraybackslash}p{0.43\tableavail}>{\raggedright\arraybackslash}p{0.29\tableavail}>{\raggedright\arraybackslash}p{0.18\tableavail}@{}}
\toprule
\textbf{Step} & \textbf{Exact command / action} & \textbf{Expected result} & \textbf{Evidence / disposition}\\
\midrule
A01 & \path{export KIWI_RUN_DIR="$PWD/docs/build/kiwi-run-001"} & Absolute new experiment location selected. Change the suffix for another run. & Record selected path; no existing evidence overwritten.\\[5pt]
A02 & \path{mkdir "$KIWI_RUN_DIR"} & Exit 0; directory created. Existing directory must cause a stop. & New empty workspace.\\[5pt]
A03 & \path{git rev-parse HEAD > "$KIWI_RUN_DIR/linkwise-commit.txt"} & One full commit ID written. Also retain the working-tree state in A04. & Source identity, not proof of a clean tree.\\[5pt]
A04 & \path{git status --porcelain=v1 > "$KIWI_RUN_DIR/working-tree.txt"} & File records every local change. A qualification baseline requires a reviewed committed state or archived source diff. & Do not silently call an edited tree the named commit.\\[5pt]
A05 & \path{cargo build --locked --manifest-path rust/Cargo.toml --features demo --bins} & Exit 0; \path{rust/target/debug/linkwise} and \path{modem-demo} available. & Build output; record Rust/toolchain versions.\\[5pt]
A06 & \path{python3 tools/test-modem.py} & Pinned reference/profile hashes match and native tests pass; exit 0. Downloads may occur. & Full test log; does not establish RF interoperability.\\[5pt]
\bottomrule
\end{tabular}
\end{adjustbox}
\captionof{table}{A: local baseline and source identity}
\end{landscape}
\begin{landscape}\thispagestyle{plain}\centering
\subsubsection{B --- Prepare and verify the 30-burst oracle}

\setlength{\tabcolsep}{5pt}\renewcommand{\arraystretch}{1.25}
\setlength{\tableavail}{\dimexpr\linewidth-6\tabcolsep\relax}
\begin{adjustbox}{max totalsize={0.96\linewidth}{0.76\textheight},center}
\begin{tabular}{@{}>{\raggedright\arraybackslash}p{0.1\tableavail}>{\raggedright\arraybackslash}p{0.43\tableavail}>{\raggedright\arraybackslash}p{0.29\tableavail}>{\raggedright\arraybackslash}p{0.18\tableavail}@{}}
\toprule
\textbf{Step} & \textbf{Exact command / action} & \textbf{Expected result} & \textbf{Evidence / disposition}\\
\midrule
B01 & Terminal A: \path{rust/target/debug/modem-demo --port 8789} & Binds localhost and prints its URL. Keep running through B03. A port conflict is a stop, not a reason to terminate an unknown process. & Local service only; no radio or Kiwi stream.\\[5pt]
B02 & Terminal B: \path{curl --fail --max-time 10 http://127.0.0.1:8789/api/config --output "$KIWI_RUN_DIR/demo-config.json"} & Exit 0; JSON advertises native levels 1--3. & Retain configuration; inspect available levels.\\[5pt]
B03 & \path{python3 tools/kiwisdr-bench.py prepare --url http://127.0.0.1:8789 --output "$KIWI_RUN_DIR/stimuli"} & Creates 30 mono 48 kHz PCM16 WAVs, \path{manifest.json}, and \path{schedule.json}. IDs B001--B030, payload prefixes K01--K30 followed by six PRBS9-derived symbols. One frame per burst. & Rust verifies the source bytes before export. The schedule records the pattern, ordering, and waveform seed.\\[5pt]
B04 & \path{python3 tools/kiwisdr-bench.py verify --manifest "$KIWI_RUN_DIR/stimuli/manifest.json" --decoder rust/target/debug/linkwise --output "$KIWI_RUN_DIR/local-verification.json"} & Exit 0 and summary \path{passed: true}, \path{files: 30}, \path{verified_frames: 30}. Every WAV hash/header and decoded level/control/payload must match its oracle. & Retain the full per-file decoder report. Any mismatch stops the campaign.\\[5pt]
\bottomrule
\end{tabular}
\end{adjustbox}
\captionof{table}{B: native 30-burst preparation and exact verification}
\end{landscape}
\begin{landscape}\thispagestyle{plain}\centering
\setlength{\tabcolsep}{5pt}\renewcommand{\arraystretch}{1.25}
\setlength{\tableavail}{\dimexpr\linewidth-6\tabcolsep\relax}
\begin{adjustbox}{max totalsize={0.96\linewidth}{0.76\textheight},center}
\begin{tabular}{@{}>{\raggedright\arraybackslash}p{0.1\tableavail}>{\raggedright\arraybackslash}p{0.43\tableavail}>{\raggedright\arraybackslash}p{0.29\tableavail}>{\raggedright\arraybackslash}p{0.18\tableavail}@{}}
\toprule
\textbf{Step} & \textbf{Exact command / action} & \textbf{Expected result} & \textbf{Evidence / disposition}\\
\midrule
B05 & \path{python3 -m json.tool "$KIWI_RUN_DIR/stimuli/schedule.json"} & Planned pre/post-roll are 30 s each; minimum interburst guard 2 s; durations are derived from actual WAV samples; capture duration is finite. & This is a proposed schedule, not actual RF timing.\\[5pt]
B06 & \path{shasum -a 256 "$KIWI_RUN_DIR/stimuli/manifest.json" "$KIWI_RUN_DIR/stimuli/schedule.json" > "$KIWI_RUN_DIR/stimulus-hashes.txt"} & Two SHA-256 records written for the frozen oracle/schedule. & Preserve these alongside per-WAV hashes in the manifest.\\[5pt]
B07 & Terminal A: press Ctrl+C. & The owned local service stops; exported evidence remains. & No receiver slots or PTT are involved.\\[5pt]
B08 & Readiness review: inspect gap detection, import conversion tests, and resource evidence required by K0. & Mark K0 complete only when all Section~\ref{md:9} conditions are evidenced. B04 alone covers clean source replay, not Kiwi continuity or resampling qualification. & STOP if missing; no RF stage may inherit a false K0 pass.\\[5pt]
\bottomrule
\end{tabular}
\end{adjustbox}
\captionof{table}{B: oracle review and K0 stop condition}
\end{landscape}
\subsubsection{C --- Provision the reviewed capture client}
The recorder is pinned to \path{4eb733e6b6147f7fbeb97ced64cdac029b202d18}. Command flags were checked against that revision's option declarations; \href{../../../rust/docs/../../docs/testing/data/kiwiclient-command-reference.json}{the source-check record} contains the URL, source hash and reviewed flags. This is source inspection, not a claim of a successful live recorder run. A missing dependency is resolved before capture; never install or change a decoder dependency mid-evaluation.

\begin{landscape}\thispagestyle{plain}\centering
\setlength{\tabcolsep}{5pt}\renewcommand{\arraystretch}{1.25}
\setlength{\tableavail}{\dimexpr\linewidth-6\tabcolsep\relax}
\begin{adjustbox}{max totalsize={0.96\linewidth}{0.76\textheight},center}
\begin{tabular}{@{}>{\raggedright\arraybackslash}p{0.1\tableavail}>{\raggedright\arraybackslash}p{0.43\tableavail}>{\raggedright\arraybackslash}p{0.29\tableavail}>{\raggedright\arraybackslash}p{0.18\tableavail}@{}}
\toprule
\textbf{Step} & \textbf{Exact command / action} & \textbf{Expected result} & \textbf{Evidence / disposition}\\
\midrule
C01 & \path{export KIWI_CLIENT_DIR="$KIWI_RUN_DIR/kiwiclient"} & New checkout path below the experiment directory. & Do not point at an unreviewed working checkout.\\[5pt]
C02 & \path{git clone --no-checkout https://github.com/jks-prv/kiwiclient.git "$KIWI_CLIENT_DIR"} & Repository fetched; no receiver connection. & Retain clone errors if acquisition fails.\\[5pt]
C03 & \path{git -C "$KIWI_CLIENT_DIR" checkout --detach 4eb733e6b6147f7fbeb97ced64cdac029b202d18} & Exact reviewed revision checked out. & Capture launcher independently checks HEAD.\\[5pt]
C04 & \path{python3 -m venv "$KIWI_RUN_DIR/capture-venv"} & Isolated capture environment created. & Record interpreter version.\\[5pt]
\bottomrule
\end{tabular}
\end{adjustbox}
\captionof{table}{C: pinned recorder checkout}
\end{landscape}
\begin{landscape}\thispagestyle{plain}\centering
\setlength{\tabcolsep}{5pt}\renewcommand{\arraystretch}{1.25}
\setlength{\tableavail}{\dimexpr\linewidth-6\tabcolsep\relax}
\begin{adjustbox}{max totalsize={0.96\linewidth}{0.76\textheight},center}
\begin{tabular}{@{}>{\raggedright\arraybackslash}p{0.1\tableavail}>{\raggedright\arraybackslash}p{0.43\tableavail}>{\raggedright\arraybackslash}p{0.29\tableavail}>{\raggedright\arraybackslash}p{0.18\tableavail}@{}}
\toprule
\textbf{Step} & \textbf{Exact command / action} & \textbf{Expected result} & \textbf{Evidence / disposition}\\
\midrule
C05 & \path{export KIWI_PY="$KIWI_RUN_DIR/capture-venv/bin/python"} & Capture interpreter is explicitly selected. & No dependence on an activated shell environment.\\[5pt]
C06 & \path{"$KIWI_PY" -m pip install numpy pypng} & Dependencies install successfully. These packages must be frozen and reviewed in C07 before testing. & Installation log; this does not qualify resampling.\\[5pt]
C07 & \path{"$KIWI_PY" -m pip freeze > "$KIWI_RUN_DIR/capture-dependencies.txt"} & Actual installed versions recorded. Freeze this environment for the campaign. & Archive with Python version and platform.\\[5pt]
C08 & \path{"$KIWI_PY" "$KIWI_CLIENT_DIR/kiwirecorder.py" --help > "$KIWI_RUN_DIR/recorder-help.txt"} & Exit 0; reviewed flags present; no receiver stream opened. & Compare flags with pinned source record.\\[5pt]
\bottomrule
\end{tabular}
\end{adjustbox}
\captionof{table}{C: capture environment and help verification}
\end{landscape}
\subsubsection{D --- Configure and preflight one receiver}
For K1, use a single cooperating candidate. For K2, repeat D and E01 in a separate shell for each agreed receiver (maximum three). R02's endpoint is an example from the reviewed registry, not a booking. Create separately reviewed profiles for R03 and R05 rather than silently reusing R02's host/gain values.

\begin{landscape}\thispagestyle{plain}\centering
\setlength{\tabcolsep}{5pt}\renewcommand{\arraystretch}{1.25}
\setlength{\tableavail}{\dimexpr\linewidth-6\tabcolsep\relax}
\begin{adjustbox}{max totalsize={0.96\linewidth}{0.76\textheight},center}
\begin{tabular}{@{}>{\raggedright\arraybackslash}p{0.1\tableavail}>{\raggedright\arraybackslash}p{0.43\tableavail}>{\raggedright\arraybackslash}p{0.29\tableavail}>{\raggedright\arraybackslash}p{0.18\tableavail}@{}}
\toprule
\textbf{Step} & \textbf{Exact command / action} & \textbf{Expected result} & \textbf{Evidence / disposition}\\
\midrule
D01 & \path{cp docs/testing/procedures/capture-profile.env.example "$KIWI_RUN_DIR/capture-profile.env"} & Profile template copied into the new run. & Original template stays unchanged.\\[5pt]
D02 & MANUAL: fill the empty profile fields from the approved station/receiver configuration, then run \path{. "$KIWI_RUN_DIR/capture-profile.env"} & Node, host/port, USB carrier kHz, passband edges Hz, fixed gain and capture duration are explicitly reviewed. K1 duration example is 120 s, subject to the agreed window. & Empty settings must stop capture. Do not guess a transmit frequency.\\[5pt]
D03 & \path{curl --fail --max-time 10 "http://$KIWI_HOST:$KIWI_PORT/status" --output "$KIWI_RUN_DIR/$KIWI_NODE_ID-status.txt"} & Metadata request returns the expected receiver name/coordinates and current status. & A listing, HTTP 200 or spare slot is not owner consent. No audio stream yet.\\[5pt]
D04 & MANUAL: complete the station readiness record in Section~\ref{md:6} and review K0, conducted or inhibited rehearsal results, and receiver terms. & Test lead records disposition before E01. Station operator authorizes the bounded first RF burst. & Stop if any required result or receiver window is missing.\\[5pt]
D05 & For K2, set \path{export KIWI_CAPTURE_S="$(python3 -c 'import json,os; print(json.load(open(os.environ["KIWI_RUN_DIR"]+"/stimuli/schedule.json"))["capture_seconds"])')"} & Duration comes from the actual prepared 30-burst schedule and fits the agreed receiver/station window. & This step precedes capture; do not extend a live session silently.\\[5pt]
\bottomrule
\end{tabular}
\end{adjustbox}
\captionof{table}{D: receiver profile and readiness}
\end{landscape}
\begin{landscape}\thispagestyle{plain}\centering
\subsubsection{E --- Capture and execute the scheduled block}

\setlength{\tabcolsep}{5pt}\renewcommand{\arraystretch}{1.25}
\setlength{\tableavail}{\dimexpr\linewidth-6\tabcolsep\relax}
\begin{adjustbox}{max totalsize={0.96\linewidth}{0.76\textheight},center}
\begin{tabular}{@{}>{\raggedright\arraybackslash}p{0.1\tableavail}>{\raggedright\arraybackslash}p{0.43\tableavail}>{\raggedright\arraybackslash}p{0.29\tableavail}>{\raggedright\arraybackslash}p{0.18\tableavail}@{}}
\toprule
\textbf{Step} & \textbf{Exact command / action} & \textbf{Expected result} & \textbf{Evidence / disposition}\\
\midrule
E01 & \path{sh docs/testing/procedures/capture-audio.sh} & One uncompressed USB audio recording is saved at its received sample rate, with \path{capture.log}. No squelch, noise blanker, resampling, exemption password or unlimited retry option is requested. & Read the actual output filename from the log. The launcher refuses an existing capture directory or wrong client commit.\\[5pt]
E02 & MANUAL: after confirming actual audio capture, wait at least 30 s before the first planned burst. & A continuous pre-roll exists; actual clock/start evidence is logged. & Network request/packet arrival time is not an RF timestamp.\\[5pt]
E03 & MANUAL, station-specific: execute the separately qualified playback/PTT procedure for the frozen WAVs and schedule. K1 sends only B001; K2 sends B001--B030 in order. & Local monitor confirms each actual transmission, level, ID and time; at least 2 s quiet between bursts. & No universal station command exists in Linkwise yet. Insert the real command/procedure revision before execution; do not transmit until then.\\[5pt]
E04 & MANUAL: preserve at least 30 s post-roll after the last actual burst. & Full last burst and trailing guards recorded before capture stops. & Record early/partial termination instead of claiming full observation.\\[5pt]
\bottomrule
\end{tabular}
\end{adjustbox}
\captionof{table}{E: bounded capture and station execution sequence}
\end{landscape}
Run E01 in the capture terminal while the station operator performs E02--E04 concurrently using the separate qualified station console. Do not wait for E01 to finish before beginning the scheduled RF block. Normal return is checked after post-roll, before preservation and restoration.

The versioned launcher contains the exact \path{kiwirecorder.py} command: USB mode, explicit \path{-L}, \path{-H} and \path{-g}, \path{--ncomp}, \path{--tlimit}, five-second connect/busy timeouts and one retry each. It applies an independent subprocess wall limit of capture duration plus 30 s. Normal completion is required; timeout/nonzero exit marks affected data ineligible pending inspection. It opens no transmitter. Inspect \href{../../../rust/docs/../../docs/testing/procedures/capture-audio.sh}{the capture script} before use; script presence or syntax checking is not live qualification.

\begin{landscape}\thispagestyle{plain}\centering
\setlength{\tabcolsep}{5pt}\renewcommand{\arraystretch}{1.25}
\setlength{\tableavail}{\dimexpr\linewidth-6\tabcolsep\relax}
\begin{adjustbox}{max totalsize={0.96\linewidth}{0.76\textheight},center}
\begin{tabular}{@{}>{\raggedright\arraybackslash}p{0.1\tableavail}>{\raggedright\arraybackslash}p{0.43\tableavail}>{\raggedright\arraybackslash}p{0.29\tableavail}>{\raggedright\arraybackslash}p{0.18\tableavail}@{}}
\toprule
\textbf{Step} & \textbf{Exact command / action} & \textbf{Expected result} & \textbf{Evidence / disposition}\\
\midrule
E05 & \path{shasum -a 256 "$KIWI_RUN_DIR/$KIWI_NODE_ID/original/"*.wav > "$KIWI_RUN_DIR/$KIWI_NODE_ID/original-hashes.txt"} & Every retained original WAV gets a hash; a missing capture is an error. & Preserve originals even when a block failed.\\[5pt]
E06 & MANUAL: inspect sequence continuity, settings, timing confidence, saturation and actual transmission evidence. & Per-burst decode/SNR/timing eligibility recorded independently of decode success. & Current raw recorder WAVs alone do not prove no sample loss. Continuity instrumentation remains a required readiness item.\\[5pt]
E07 & MANUAL: verify all owned capture processes ended and receiver slots were released; execute station restoration. & No residual owned audio stream or PTT; terminal result logged. & Timeout/fault requires reconciliation, not automatic RF retry.\\[5pt]
\bottomrule
\end{tabular}
\end{adjustbox}
\captionof{table}{E: duration, evidence preservation and restoration}
\end{landscape}
\subsubsection{F --- Import, decode, compare, and report}
Install and qualify FFmpeg with libsoxr before this stage. It is an additional analysis dependency, not currently bundled with Linkwise. Record its version and build configuration. Before F03, assign the actual capture path with \path{export KIWI_AUDIO_WAV="/absolute/path/from-capture-log.wav"} and the reviewed integer \path{KIWI_START_SAMPLE} and \path{KIWI_END_SAMPLE} in the analysis shell. The shown path is a required substitution, not an existing fixture.

The commands below are for \textbf{demodulated mono audio only}; never downmix an IQ recording as though it were audio. \path{KIWI_AUDIO_WAV} is the actual capture filename from E01. \path{KIWI_START_SAMPLE} and \path{KIWI_END_SAMPLE} are reviewed start and end indices at the original sample rate for a complete continuous eligible interval, including background and guards. For a fully continuous capture, use zero and the original sample count. Use source sample indices rather than network arrival times. Conversion preserves sample order.

\begin{landscape}\thispagestyle{plain}\centering
\setlength{\tabcolsep}{5pt}\renewcommand{\arraystretch}{1.25}
\setlength{\tableavail}{\dimexpr\linewidth-6\tabcolsep\relax}
\begin{adjustbox}{max totalsize={0.96\linewidth}{0.76\textheight},center}
\begin{tabular}{@{}>{\raggedright\arraybackslash}p{0.1\tableavail}>{\raggedright\arraybackslash}p{0.43\tableavail}>{\raggedright\arraybackslash}p{0.29\tableavail}>{\raggedright\arraybackslash}p{0.18\tableavail}@{}}
\toprule
\textbf{Step} & \textbf{Exact command / action} & \textbf{Expected result} & \textbf{Evidence / disposition}\\
\midrule
F01 & \path{ffmpeg -version > "$KIWI_RUN_DIR/ffmpeg-version.txt"} & Installed analysis version recorded. & Missing executable blocks import.\\[5pt]
F02 & \path{ffmpeg -buildconf > "$KIWI_RUN_DIR/ffmpeg-build.txt" 2>&1} & Reviewed build supports libsoxr. & Do not silently fall back to another resampler.\\[5pt]
F03 & \path{ffprobe -v error -show_entries stream=sample_rate,channels,sample_fmt -of json "$KIWI_AUDIO_WAV" > "$KIWI_RUN_DIR/$KIWI_NODE_ID/input-format.json"} & Original received rate recorded; exactly one demodulated audio channel confirmed. & Reject IQ/stereo input for this audio-only procedure.\\[5pt]
\bottomrule
\end{tabular}
\end{adjustbox}
\captionof{table}{F: analysis dependency and audio format checks}
\end{landscape}
\begin{landscape}\thispagestyle{plain}\centering
\setlength{\tabcolsep}{5pt}\renewcommand{\arraystretch}{1.25}
\setlength{\tableavail}{\dimexpr\linewidth-6\tabcolsep\relax}
\begin{adjustbox}{max totalsize={0.96\linewidth}{0.76\textheight},center}
\begin{tabular}{@{}>{\raggedright\arraybackslash}p{0.1\tableavail}>{\raggedright\arraybackslash}p{0.43\tableavail}>{\raggedright\arraybackslash}p{0.29\tableavail}>{\raggedright\arraybackslash}p{0.18\tableavail}@{}}
\toprule
\textbf{Step} & \textbf{Exact command / action} & \textbf{Expected result} & \textbf{Evidence / disposition}\\
\midrule
F04 & \path{mkdir "$KIWI_RUN_DIR/$KIWI_NODE_ID/derived"} & New derivative directory created. & Existing evidence directory must not be overwritten.\\[5pt]
F05 & \path{ffmpeg -nostdin -n -i "$KIWI_AUDIO_WAV" -af "atrim=start_sample=$KIWI_START_SAMPLE:end_sample=$KIWI_END_SAMPLE,asetpts=PTS-STARTPTS,aresample=48000:resampler=soxr:precision=28" -ac 1 -c:a pcm_s16le "$KIWI_RUN_DIR/$KIWI_NODE_ID/derived/capture-48k.wav"} & New 48 kHz mono PCM16 WAV for the continuous interval. Original stays unchanged. & Save the interval mapping and filter settings. Preserve noise and guards.\\[5pt]
F06 & \path{rust/target/debug/linkwise decode-wav "$KIWI_RUN_DIR/$KIWI_NODE_ID/derived/capture-48k.wav" > "$KIWI_RUN_DIR/$KIWI_NODE_ID/decode.json"} & Valid JSON with observed frames, or an explicit decoder error retained. Zero recovered frames is a valid RF/decode miss when the observation is eligible. & Compare count, order, and bytes in F07. Preserve truncated energy candidate counts.\\[5pt]
\bottomrule
\end{tabular}
\end{adjustbox}
\captionof{table}{F: declared audio conversion and Rust decode}
\end{landscape}
\begin{landscape}\thispagestyle{plain}\centering
\setlength{\tabcolsep}{5pt}\renewcommand{\arraystretch}{1.25}
\setlength{\tableavail}{\dimexpr\linewidth-6\tabcolsep\relax}
\begin{adjustbox}{max totalsize={0.96\linewidth}{0.76\textheight},center}
\begin{tabular}{@{}>{\raggedright\arraybackslash}p{0.1\tableavail}>{\raggedright\arraybackslash}p{0.43\tableavail}>{\raggedright\arraybackslash}p{0.29\tableavail}>{\raggedright\arraybackslash}p{0.18\tableavail}@{}}
\toprule
\textbf{Step} & \textbf{Exact command / action} & \textbf{Expected result} & \textbf{Evidence / disposition}\\
\midrule
F07 & MANUAL: map each decoded sample range to its original capture interval and confirmed transmit window. Compare level, control and full payload with \path{stimuli/manifest.json}. & Exact match counts that expected identity once. Missing expected frame counts as a miss; duplicates/conflicts recorded separately. & Do not use the local-source \path{kiwisdr-bench.py verify} as an RF scorer; it verifies source files, not capture associations.\\[5pt]
F08 & MANUAL: reconcile every planned receiver/burst with TX confirmation, eligibility, exact match/miss and exclusion reason. & Denominators include eligible undetected bursts; infrastructure failures remain visible separately. & Automated RF association and aggregate reporting are still backlog items.\\[5pt]
F09 & MANUAL: repeat F05--F08 on a frozen corpus with a separately built candidate decoder and identical conversion. & Paired gains/losses, mismatches and resource results retained by receiver, level and session. & No improvement claim without both versions and the frozen oracle.\\[5pt]
\bottomrule
\end{tabular}
\end{adjustbox}
\captionof{table}{F: exact RF comparison and report disposition}
\end{landscape}
The FFmpeg syntax is based on its official \href{https://ffmpeg.org/ffmpeg-resampler.html}{resampler} and \href{https://ffmpeg.org/ffmpeg-filters.html#atrim}{atrim filter} documentation. This procedure must be qualified on known audio before RF evidence is scored; source-reviewed commands are not a completed conversion test.

For K0, use the approved FFmpeg build to derive a mono PCM16 12 kHz copy of the clean B001 source with \path{-af aresample=12000:resampler=soxr:precision=28}. Use F03--F07 on the entire clean copy, with the sample range starting at zero. Require the exact expected B001 frame and the documented 12-to-48 kHz sample mapping. Repeat with the same input and compare the derived WAV hashes and frame results. Preserve both sample counts and the FFmpeg build configuration.

For the 50 ms control, represent the 12 kHz source as 600-sample numbered blocks with their original offsets. Remove one block fully inside the known burst and retain the original block indices in the fixture sidecar. Run the same continuity check on the clean and gapped sidecars. Require a contiguous clean result and a reported 600-sample discontinuity in the gapped result. Mark the affected observation ineligible before conversion or decoding. This synthetic control verifies the check against known missing samples. K0 also requires showing how the pinned recorder produces equivalent continuity evidence for a real capture.

\subsubsection{Recorded validation of this procedure draft}
Only the file/loopback path and fail-fast controls are exercised while drafting. Record actual local results separately from expected entries above. Capture-client option inspection, shell syntax checks, source preparation and local exact-byte replay do not demonstrate an RF experiment, continuity-qualified Kiwi recording, libsoxr conversion, or an operational transmit chain.

Local draft-validation result (9 November 2026): all 30 generated WAVs decoded with exact frame agreement; the planned capture length was 277 s for this build. Wrong-payload and wrong-hash oracles returned failure, and a nonempty destination was rejected. The capture launcher passed shell syntax checks; missing profile and wrong-checkout guards were tested without opening a receiver. No live capture, radio transmission, or FFmpeg conversion was performed in these checks.

Revision 0.9 local validation (12 November 2026): the pinned modem tests passed. The updated helper generated 30 PRBS9-derived WAVs and the Rust WAV decoder recovered all 30 expected frames. An independent payload check confirmed the register phases and encoded symbols. A second preparation produced identical WAVs, manifests, and schedule. The scheduled capture length was 277 s. The available FFmpeg binary lacked libsoxr, so the K0 conversion check, recorder continuity evidence, and station rehearsal remain pending.


\appendix

\section{Implementation Backlog and Existing Commands}
\label{tp:backlog}

\begin{enumerate}
\item Use \path{tools/kiwisdr-bench.py prepare} and \path{verify} for 30 uniquely identified native source bursts and exact local replay. Qualify the remaining import and continuity tooling before treating K0 as complete.
\item The versioned capture launcher now bounds a pinned recorder invocation; qualify it with cooperating receivers and add verified continuity metadata. Preserve originals at their received rate. See the exact sequence in section 14.
\item Add deterministic audio/IQ import and a manifest validator, including gap and sample-position handling. Neither component currently exists in Linkwise.
\item Add a comparator/report harness around the existing Rust WAV CLI. Add missing acquisition/soft-decision metrics only as separately tested instrumentation.
\item Complete K0 before arranging K1 receiver access and a qualified RF transmission.
\end{enumerate}
Existing commands from the repository root:

\begin{lstlisting}
# Existing external fixture/profile validation; downloads pinned inputs.
python3 tools/test-modem.py

# Decode an already converted 48 kHz mono analysis WAV.
cargo run --locked --manifest-path rust/Cargo.toml --bin linkwise -- \
  decode-wav /path/to/derived-48k-mono.wav
\end{lstlisting}
For optional level-4 replay, set \path{LINKWISE_BW500_PROFILE} to the verified profile as described in the modem guide. Capture wrappers, comparison reports, RF playback, and global receiver selection are backlog items, not commands supplied by this draft.


\section{Template Lineage and Tailoring}
\label{app:tailoring}

The crosswalk below refers to section numbers in the companion Markdown draft, not this edition.

Use the repository's \href{https://github.com/8r4n/acq-latex-templates/blob/fef160e37542fd7bc8e62de81404a83b8634a396/templates/engineering/verification/test-plan/main.tex}{engineering verification Test Plan template} as the structural basis for this plan. Inspected source revision: \path{fef160e37542fd7bc8e62de81404a83b8634a396}, 9 November 2026. The template identifies DI-NDTI-80808 as its source; this draft does not claim independent verification of that document or contractual compliance.

This is a system/modem experiment involving software, an RF transmit chain, remote instruments, capture integrity, and data analysis. The broader test-plan template fits better than the \href{https://github.com/8r4n/acq-latex-templates/blob/fef160e37542fd7bc8e62de81404a83b8634a396/templates/engineering/software/stp/main.tex}{Software Test Plan}, which is oriented toward software qualification against software requirements. A later software qualification plan can cover the decoder/import/comparator independently without treating this exploratory RF pilot as formal qualification.

Retain the template's engineering organization and replace acquisition-program metadata with Linkwise document ID, revision, author/maintainer, test lead, configuration baseline, and review status. Do not copy placeholder signatures, agency roles, acquisition milestones, or default distribution markings into the project. This Markdown draft remains the maintained content source. The generated, template-structured LaTeX edition is \href{../../../rust/docs/../../docs/testing/kiwisdr-modem-test-plan.tex}{kiwisdr-modem-test-plan.tex}; run \path{python3 tools/build-kiwisdr-test-plan.py --build} from the repository root.

\begin{landscape}\thispagestyle{plain}\centering
\setlength{\tabcolsep}{5pt}\renewcommand{\arraystretch}{1.25}
\setlength{\tableavail}{\dimexpr\linewidth-2\tabcolsep\relax}
\begin{adjustbox}{max totalsize={0.96\linewidth}{0.76\textheight},center}
\begin{tabular}{@{}>{\raggedright\arraybackslash}p{0.29\tableavail}>{\raggedright\arraybackslash}p{0.71\tableavail}@{}}
\toprule
\textbf{Selected template section} & \textbf{Content in this draft / tailoring}\\
\midrule
Scope & Sections 1--2: objectives, boundaries, implementation baseline\\[5pt]
Referenced Documents & Section 12 and the immutable template reference here; freeze all execution dependencies before a campaign\\[5pt]
Test Items & Sections 2, 5, and 10: transmitter waveforms, decoder commits, external profile, capture/conversion pipeline, as-tested manifests\\[5pt]
Test Approach & Sections 3--9 plus the objective records below\\[5pt]
Resources, Facilities, and Equipment & Section 4 and resource/status table below\\[5pt]
Safety & Station readiness, abort/restore rules below; campaign-specific hazards remain to be assessed before RF\\[5pt]
Schedule & K0--K5 dependency sequence; calendar windows remain unscheduled until resources and entry criteria are confirmed\\[5pt]
Data Management and Reporting & Sections 8 and 10: immutable originals, derived artifacts, scoring, exclusions, and final report\\[5pt]
Submission-readiness appendix & Use the readiness review below, adapted to a project experiment rather than a contractual submission\\[5pt]
\bottomrule
\end{tabular}
\end{adjustbox}
\captionof{table}{Template crosswalk into the maintained Markdown draft}
\end{landscape}

\clearpage
\section{Campaign-readiness Checklist}
\label{app:checklist}

\begin{itemize}
\item Test lead and station operator assigned; no placeholder approval signatures treated as authorization.
\item Matrix, waveform hashes, expected bytes, configuration and plan/procedure revisions frozen.
\item K0 tooling evidence complete, including gap detection and exact-byte import checks.
\item Receiver windows/settings/terms confirmed; station transmit and stop path qualified.
\item Recording duration, airtime, resource limits, storage and restoration actions verified.
\item Metric eligibility and exclusions declared; no decoding success used to hide capture failures.
\item Report links each objective and verdict to as-tested configuration and raw/reduced evidence.
\end{itemize}
This is a planned readiness record. No item is represented as complete by inclusion in this checklist.


\end{document}
